\documentclass[12pt]{article}

%% ── Packages ──────────────────────────────────────────────────────────────
\usepackage[margin=1.2in]{geometry}
\usepackage{amsmath,amssymb,amsthm}
\usepackage{mathtools}
\usepackage{hyperref}
\usepackage{microtype}
\usepackage{booktabs}
\usepackage{array}
\usepackage{enumitem}
\usepackage{xcolor}
\usepackage{tcolorbox}
\usepackage{setspace}
\usepackage[numbers,sort&compress]{natbib}

%% ── Theorem environments ──────────────────────────────────────────────────
\newtheorem{theorem}{Theorem}[section]
\newtheorem{lemma}[theorem]{Lemma}
\newtheorem{corollary}[theorem]{Corollary}
\newtheorem{definition}[theorem]{Definition}
\newtheorem{proposition}[theorem]{Proposition}
\theoremstyle{remark}
\newtheorem{remark}[theorem]{Remark}

%% ── Custom commands ───────────────────────────────────────────────────────
\newcommand{\Z}{\mathbb{Z}}
\newcommand{\R}{\mathbb{R}}
\newcommand{\N}{\mathbb{N}}
\newcommand{\abs}[1]{\lvert #1 \rvert}
\newcommand{\norm}[1]{\lVert #1 \rVert}
\newcommand{\ceil}[1]{\lceil #1 \rceil}
\newcommand{\floor}[1]{\lfloor #1 \rfloor}
\DeclareMathOperator{\lcm}{lcm}

%% ── Provenance box ────────────────────────────────────────────────────────
\tcbuselibrary{breakable,skins}
\newtcolorbox{provenancebox}{
  breakable,
  colback=gray!8,
  colframe=gray!50,
  title={\textbf{Provenance and Status}},
  fonttitle=\small\bfseries,
  fontupper=\small,
  left=6pt, right=6pt, top=4pt, bottom=4pt
}

%% ── Metadata ──────────────────────────────────────────────────────────────
\title{%
  Exploring Constructive Approaches to the Umans--Wang\\
  $(\alpha,\beta)$-Divisor Conjecture:\\[4pt]
  \large A Survey of Construction Strategies, Empirical Obstructions,\\
  and Lattice-Based Optimality Results%
}
\author{Prepared in collaboration with Claude (Anthropic)}
\date{2026 --- Revised following peer review}

%% ═══════════════════════════════════════════════════════════════════════════
\begin{document}
\maketitle
\thispagestyle{empty}

%% ── Provenance ────────────────────────────────────────────────────────────
\begin{provenancebox}
This document was produced through an extended, iterative conversation between a
human collaborator and Claude, an AI model developed by Anthropic.  Real errors
were introduced and subsequently caught during the investigation: a sign error in
an early continued-fraction computation; a degenerate zero-valued covering term
that initially passed verification spuriously; a random-seed bug making nominally
independent trials identical; and a floating-point precision failure that
produced confident but wrong numerical results before diagnosis.  These are
disclosed because they are relevant to how the results should be weighed.

This work has \textbf{not} been peer-reviewed, has not been checked by a domain
expert in number theory or computational complexity, and has not been submitted to
any venue.  The three theorems are carefully checked but not independently
verified.  The non-degeneracy lemma (Lemma~\ref{lem:nondegen}) rests on one
cited classical fact (Lemma~\ref{lem:bounded-basis}) rather than a fully
self-contained derivation.  The numerical results in Section~\ref{sec:star}
characterise the empirical performance of our specific lattice-reduction
algorithm relative to a target that Theorem~\ref{thm:fullline} already
establishes unconditionally; they should not be read as evidence that practical
performance converges to $3/2$ independently of that theorem.  The literature
search supporting novelty claims was targeted, not exhaustive.
\end{provenancebox}

\vspace{8pt}
\begin{abstract}
Umans and Wang~\cite{umans2025} introduced the $(\alpha,\beta)$-Divisor
Conjecture: whether there exist sets $S, T \subseteq \N$, each of cardinality at
most $n^{\beta}$, with elements of magnitude at most $\exp(n^{\alpha})$, such
that every integer $i \le n$ divides some difference $s - t$ with $s \in S$,
$t \in T$.  A positive resolution with $\alpha, \beta < 1/2$ would yield the
first improvement in over a decade to the best known algorithms for polynomial
factorisation over finite fields and for deterministic integer factorisation.

We document a systematic exploration of construction strategies spanning
elementary number theory, the probabilistic method, continued fractions and
Diophantine approximation, discrete logarithms, the geometry of numbers, and
classical addition-chain theory.  We report three new, fully proved results:
\begin{enumerate}[label=(\roman*)]
  \item an exact characterisation of the optimal trade-off achievable by a
    natural family of \emph{hypercube lattice} constructions
    (Theorem~\ref{thm:hypercube});
  \item a strictly stronger \emph{full-line achievability} theorem showing a
    related star/chain construction reaches the trivial bound $\alpha+\beta=1$
    continuously (Theorem~\ref{thm:fullline});
  \item a unifying theorem proving the star/chain structure is the global optimum
    of the entire two-parameter family of blocked hypercube--star hybrids, via
    the one-line inequality $2^d - 1 \ge d$ (Theorem~\ref{thm:global}).
\end{enumerate}
We construct explicit, verified witnesses approaching the exponent of the best
previously known algorithm~\cite{kedlaya2011}, and examine carefully what those
experiments do and do not establish.  The $(\alpha,\beta)$-Divisor Conjecture
remains open.
\end{abstract}

\tableofcontents
\newpage

%% ═══════════════════════════════════════════════════════════════════════════
\section{The Umans--Wang Conjecture}

In November 2025, Umans and Wang~\cite{umans2025} posed a number-theoretic
conjecture with a striking algorithmic consequence.  The current best algorithm
for factoring polynomials over finite fields runs in time roughly $n^{3/2}$, a
barrier standing since Kedlaya--Umans~\cite{kedlaya2011}; the conjecture, if
true with the right parameters, would push this exponent to $n^{4/3}$.

\begin{definition}[$(\alpha,\beta)$-covering pair]
For $n \in \N$ and $\alpha,\beta > 0$, a pair $(S,T)$ of sets of positive
integers is an \emph{$(\alpha,\beta)$-covering pair for $n$} if $|S|, |T| \le
n^{\beta}$, every element of $S \cup T$ is at most $\exp(n^{\alpha})$, and for
every $i \in \{1,\ldots,n\}$ there exist $s \in S$, $t \in T$ with $i \mid s -
t$.
\end{definition}

\begin{definition}[$n$-divisor property~{\cite[Def.~3.1]{umans2025}}]
A set $A \subseteq \N$ satisfies the \emph{$n$-divisor property} if for every
$i \in \{1,\ldots,n\}$ there exists $a \in A$ with $i \mid a$.
\end{definition}

The \emph{Strong $(\alpha,\beta)$-Divisor Conjecture} asks whether covering
pairs exist with $\alpha,\beta < 1/2$.  Two trivial constructions span the easy
ends: $S = \{1,\ldots,n\}$, $T = \{0\}$ gives $\alpha \approx 0$, $\beta = 1$;
$S = \{\lcm(1,\ldots,n)\}$, $T = \{0\}$ gives $\beta = 0$, $\alpha \approx 1$.
Both sit on the line $\alpha + \beta = 1$.

%% ═══════════════════════════════════════════════════════════════════════════
\section{Elementary Constructions and the Trivial Floor}

Bucketing constructions --- partitioning $[1,n]$ into $m$ groups and forming
$S$ from group LCMs --- achieve $\beta = \log(m)/\log(n)$ and $\alpha = 1 -
\beta + o(1)$, tracing the trivial line exactly.  Arithmetic-progression
variants, motivated by Proposition~3.4 of~\cite{umans2025} (which reduces the
conjecture to a single-AP question), gave the same trade-off in every
configuration tested.  Greedy combinatorial search over smooth-number candidates
similarly landed at or above $\alpha + \beta \approx 1$ across all problem
sizes, consistent with the combinatorial impossibility result of
Lemma~6.3 in~\cite{umans2025}.

%% ═══════════════════════════════════════════════════════════════════════════
\section{An Exact Construction via the Chinese Remainder Theorem}

Lemma~6.2 of~\cite{umans2025} gives an exact algebraic characterisation: an
arithmetic progression covers every prime $p \le n$ if and only if there is a
factorisation of the primorial $U = \prod_{p \le n} p$ together with integers
$b, c$ satisfying
\[
  c \cdot S + b \equiv 0 \pmod{U},
\]
where $S$ is a sum of Chinese Remainder idempotents determined by the
factorisation.  For fixed factorisation, finding the smallest $(b,c)$ pair
reduces to continued-fraction approximation of $S/U$.  After correcting an
initial sign error ($b \equiv +cS$ instead of $-cS$) and a degenerate zero-term
bug (a covering term equal to $0$, trivially divisible by every prime but
mathematically meaningless), the implementation produced verified covering
progressions with magnitudes well below the generic $\sqrt{U}$ baseline.

A scaling study from $n = 10$ to $n = 150$ showed real but logarithmic
magnitude savings consistent with the metric theory of continued fractions
(Khinchin~\cite{khinchin1964}, Postnikov~\cite{postnikov1966}): the typical
partial quotient in the continued fraction of a ``generic'' number grows like
$O(\log k)$, so the savings are real but do not change the leading-order
exponent.

%% ═══════════════════════════════════════════════════════════════════════════
\section{The Discrete-Logarithm Sieve}

For a fixed base $g$ and prime $p$, one has $p \mid g^i - 1$ if and only if
$\mathrm{ord}_p(g) \mid i$.  Taking $S = \{g, g^2, \ldots, g^L\}$ and $T =
\{1\}$ therefore covers any prime whose multiplicative order under $g$ divides
some $i \le L$.  Direct computation of orders across primes up to several
hundred showed coverage growing essentially linearly with $L$, with no anomalous
smoothness for any fixed small base --- consistent with known results on the
average multiplicative order of a fixed integer mod $p$.  A hybrid construction
pre-filtering smooth-order primes before handing the residual to the
continued-fraction search was implemented and found to have a strictly worse
combined $(\alpha,\beta)$ trade-off than the continued-fraction method alone in
every configuration tested.

%% ═══════════════════════════════════════════════════════════════════════════
\section{The Probabilistic Method}

\subsection{First moment}

For $S, T$ chosen as uniform random subsets of size $K$ in a large range, the
expected number of uncovered divisors $i \le n$ is $\sum_{i=1}^{n}
\exp(-K^2/i)$.  A union bound shows this expectation falls below $1$ once $K
\gtrsim \sqrt{n \log n}$, giving a threshold $\beta \to 1/2$ as $n \to \infty$.
Numerical verification across $n \in \{50, 100, 500, 1000\}$ confirmed
convergence of the implied $\beta$ to $1/2$ exactly, independently recovering
the trivial construction's floor via a purely probabilistic argument.

\subsection{Structured randomness}

Sampling $S$ and $T$ as random arithmetic progressions rather than uniform sets
was tested directly and found to perform no better, and measurably worse, than
uniform sampling.  The first-moment calculation depends only on marginal
probabilities that a single pair $(s,t)$ satisfies divisibility, which
structured sampling does not improve; AP-structured differences are less varied
than uniform ones, which can only hurt coverage.

\subsection{Dispersion method}

An attempt to apply Linnik's dispersion method~\cite{linnik1963} (the
second-moment escalation used to solve binary additive problems such as the
Titchmarsh divisor problem) found that the events ``$i$ is uncovered'' and
``$j$ is uncovered'' are \emph{positively} correlated: failing to cover one
divisor increases the probability of failing to cover another.  This is the
wrong sign for Janson-inequality refinements to help.  We computed this for one
representative pair $(i,j) = (40, 60)$ analytically and by simulation; we did
not compute the aggregate dependency term $\Delta = \sum_{i,j} \mathrm{Cov}(X_i,
X_j)$ required for a fully rigorous conclusion, so this should be read as a
plausible mechanistic explanation rather than a proof.

%% ═══════════════════════════════════════════════════════════════════════════
\section{The Alteration Method}

Deliberately undershooting the random-construction threshold by roughly $30\%$
and patching each leftover uncovered divisor with one targeted extra element
achieved a real $20$--$25\%$ reduction in total construction size relative to
forcing full coverage outright.  However, composing this with the exact
CRT/continued-fraction construction of Section~3 revealed a structural mismatch:
the continued-fraction construction succeeds with certainty for any valid
factorisation (by B\'ezout's identity), leaving no probabilistic near-miss for
the alteration method to patch.  Across hundreds of trials at $n$ up to $800$,
the construction achieved complete coverage in every instance.

%% ═══════════════════════════════════════════════════════════════════════════
\section{Generalised Progressions and the Lattice Method}

Conjecture~3.3 of~\cite{umans2025} requires $S$ and $T$ to be sums of several
arithmetic progressions.  Testing the two-progression case via Lemma~6.1's
combination formula found it strictly worse than the single-progression
baseline: the count overhead of the combination formula was not offset by a
sufficient reduction in magnitude.

Reformulating as a joint optimisation revealed the problem reduces to a lattice
question.  With $k$ independent common differences $b_1,\ldots,b_k$ and shared
offset $B$, the covering condition collapses via CRT idempotents to
\[
  B + \sum_{i=1}^{k} b_i W_i \equiv 0 \pmod{U},
\]
where $W_i \in \Z/U\Z$ are constants determined by the prime partition.  This
is a rank-$(k+1)$ lattice of covolume $U$, and Minkowski's First Theorem
guarantees a nonzero point with all coordinates bounded by
$C \cdot U^{1/(k+1)}$.

%% ═══════════════════════════════════════════════════════════════════════════
\section{A Proved Result: Optimality of Hypercube Lattice Constructions}
\label{sec:hypercube}

\subsection{Theorem 1}

\begin{theorem}[Optimality of hypercube constructions]
\label{thm:hypercube}
For the hypercube construction with $k$ binary dimensions indexed by $\varepsilon
\in \{0,1\}^k$ (i.e.\ $S = \{B + \sum_i \varepsilon_i b_i : \varepsilon \in
\{0,1\}^k\}$), the rank-$(k+1)$ lattice of covolume $U$ yields trade-off
\[
  \alpha + \beta
  = 1 + \frac{k \ln 2 - \ln(k+1)}{\ln n} + o(1).
\]
This is minimised over non-negative integers $k$ at $k \in \{0,1\}$, where it
equals $1$ exactly.  For all integers $k \ge 2$ it is strictly greater than $1$.
\end{theorem}

\begin{proof}
(a)~\textbf{Covolume.}  The basis matrix is lower-triangular with $U$ in
position $(0,0)$ and ones on the remaining diagonal, so $\det = U$ exactly,
independent of the partition and the $W_i$ values.

(b)~\textbf{Coverage.}  For prime $p$ assigned to class $\varepsilon^{(p)}$, the
idempotent property gives $W_i \equiv \varepsilon_i^{(p)} \pmod{p}$ for each
$i$, so the lattice relation $B + \sum_i b_i W_i \equiv 0 \pmod{U}$ implies $B
+ \sum_i \varepsilon_i^{(p)} b_i \equiv 0 \pmod{p}$, i.e.\ the term of $S$
indexed by $\varepsilon^{(p)}$ is divisible by $p$.

(c)~\textbf{Minkowski bound.}  By Minkowski's First Theorem, the lattice
contains a nonzero point with $\ell^\infty$-norm at most $C(d,N) \cdot
U^{1/(k+1)}$ avoiding $N = 2^k$ forbidden hyperplanes (one per coordinate
choice); the non-degeneracy of this point is guaranteed by
Lemma~\ref{lem:nondegen} below.

(d)~\textbf{Trade-off.}  By the prime number theorem $\log U = (1+o(1))n$.
Hence $\alpha = 1 - \ln(k+1)/\ln n + o(1)$ and $\beta = k \ln 2 / \ln n$.
Setting $f(k) = k \ln 2 - \ln(k+1)$: $f(0) = f(1) = 0$ exactly; $f'(k) = \ln 2
- 1/(k+1) = 0$ at $k+1 = 1/\ln 2 \approx 1.44$; $f$ is convex and positive for
all integers $k \ge 2$ (e.g.\ $f(2) = 2\ln 2 - \ln 3 \approx 0.288 > 0$). \qed
\end{proof}

\subsection{Non-degeneracy lemma}
\label{sec:nondegen}

The step ``there exists a nonzero lattice point avoiding the $N$ forbidden
hyperplanes'' requires proof.

\begin{lemma}[Bounded basis]
\label{lem:bounded-basis}
Every rank-$d$ lattice $L$ of covolume $U$ admits a basis $v_1,\ldots,v_d$
with $\norm{v_i}_\infty \le C_1(d) \cdot U^{1/d}$ for all $i$, where $C_1(d)$
depends only on $d$.
\end{lemma}

This is standard: Minkowski's Second Theorem bounds $\lambda_1 \cdots \lambda_d
\le U$; classical reduction theory (see Cassels~\cite{cassels1959}, Ch.~VIII)
gives an actual basis bounded in terms of the successive minima.  We cite this
without re-deriving the sharp constant; a crude bound $\norm{v_i}_\infty \le
d \cdot U$ suffices for our application.

\begin{lemma}[Non-degeneracy]
\label{lem:nondegen}
Let $L$ be a rank-$d$ lattice of covolume $U$ ($d \ge 2$), and let
$H_1,\ldots,H_N \subset \R^d$ be hyperplanes through the origin.  Then there
exists a nonzero $v \in L$ with $\norm{v}_\infty \le N \cdot d \cdot C_1(d)
\cdot U^{1/d}$ lying on none of $H_1,\ldots,H_N$.
\end{lemma}

\begin{proof}
Fix a basis $v_1,\ldots,v_d$ from Lemma~\ref{lem:bounded-basis}.  For $c =
(c_1,\ldots,c_d) \in \{0,\ldots,N\}^d$ set $v(c) = \sum_i c_i v_i$.  There
are $(N+1)^d - 1$ nonzero such combinations.

Fix hyperplane $H_j$ with nonzero functional $\ell_j$.  Since $v_1,\ldots,v_d$
are linearly independent and $\ell_j \ne 0$, there exists $m$ with $\ell_j(v_m)
\ne 0$.  Fixing the remaining $d-1$ coefficients, $v(c) \in H_j$ constrains
$c_m$ uniquely, so $H_j$ captures at most $(N+1)^{d-1}$ tuples.

By the union bound, the number of nonzero tuples landing on at least one $H_j$
is at most $N(N+1)^{d-1}$.  Good tuples number at least
\[
  (N+1)^d - 1 - N(N+1)^{d-1} = (N+1)^{d-1} - 1 \ge 1
  \quad (d \ge 2,\; N \ge 1).
\]
Any such $c^*$ gives $v = v(c^*) \ne 0$ with $\norm{v}_\infty \le N \cdot
\max_i \norm{v_i}_\infty \le N \cdot d \cdot C_1(d) \cdot U^{1/d}$.  \qed
\end{proof}

\begin{remark}[Well-roundedness not required]
The proof requires only that $\ell_j$ cannot vanish on \emph{every} basis vector
simultaneously (guaranteed by linear independence), not that each basis vector
individually avoids each hyperplane.  If a short basis vector $v_m$ happens to
lie on $H_j$, its coefficient drops out of the linear equation for $H_j$,
leaving the bound enforced through whichever other basis vector has nonzero
$\ell_j(v_i)$.  We verified this concretely on an adversarial example in which
two of three basis vectors lay exactly on the same hyperplane: the count of bad
coefficient tuples was exactly $8$ out of $26$ nonzero candidates, consistent
with (in fact equal to) the bound $(N+1)^{d-1} - 1 = 8$ for $N=2$, $d=3$.  The
lattice need not be well-rounded.
\end{remark}

\begin{remark}[Hermite's constant at $k=3$ and $k=7$]
The Minkowski bound $C_1(d) \cdot U^{1/d}$ can be sharpened using Hermite's
constant $\gamma_d$: $\gamma_4 = \sqrt{2}$ (Korkine--Zolotareff~\cite{kz1873},
achieved by the $D_4$ lattice) and $\gamma_8 = 2$ (Viazovska~\cite{viazovska2017},
$E_8$ lattice).  Substituting these exact constants at $k=3$ or $k=7$ sharpens
the leading constant $C(k)$ without affecting the asymptotic content of
Theorems~\ref{thm:hypercube}--\ref{thm:global}.
\end{remark}

%% ═══════════════════════════════════════════════════════════════════════════
\section{A Flexible Construction: Star Chains}
\label{sec:star}

\subsection{Construction}

The \emph{star/chain} construction uses $k$ generators each combining only with
a shared base $B$:
\[
  S = \{B,\; B + b_1,\; B + b_2,\; \ldots,\; B + b_k\}, \quad T = \{0\}.
\]
This gives $k+1$ terms (linear in $k$) versus $2^k$ for the hypercube, while
the underlying lattice has the same rank $k+1$ and covolume $U$.

\begin{theorem}[Full-line achievability]
\label{thm:fullline}
Fix $\beta \in (0,1)$.  For the star/chain construction with $k = \lceil n^\beta
\rceil - 1$ generators,
\[
  \beta_{\mathrm{achieved}} = \beta + o(1), \qquad
  \alpha_{\mathrm{achieved}} = 1 - \beta + o(1)
  \quad \text{as } n \to \infty,
\]
so $\alpha + \beta = 1 + o(1)$ for every fixed $\beta \in (0,1)$.
\end{theorem}

\begin{proof}
Covolume is $U$ (same argument as Theorem~\ref{thm:hypercube}).  With $k+1$
forbidden hyperplanes and $d = k+1$, Lemma~\ref{lem:nondegen} gives a valid
covering construction of magnitude at most $C(k) \cdot U^{1/(k+1)}$ with
$C(k) = O(k)$.  Then $\log(\mathrm{magnitude}) \le \log C(k) + n(1+o(1))/(k+1) =
O(\log n) + \Theta(n^{1-\beta})$.  For fixed $\beta < 1$, $n^{1-\beta}$
dominates $\log n$ as $n \to \infty$, giving $\alpha = 1 - \beta + o(1)$.

\textbf{Caveat:} convergence is pointwise in $\beta$; the rate slows as $\beta
\to 1$, since the ratio $\log C(k)/n^{1-\beta}$ shrinks only as a power of
$1/n$ for $\beta < 1$ but can remain $O(1)$ for $\beta$ close to $1$ at any
finite $n$.  \qed
\end{proof}

\subsection{Theorem 3: Global optimality of the chain structure}

\begin{definition}[Blocked hypercube--star hybrid]
For positive integers $m, d$ with $k = md$, the $(m,d)$-hybrid uses $m$ blocks
of $d$ generators each, with full $d$-dimensional hypercube structure
\emph{within} each block and star structure \emph{between} blocks.  This has
lattice rank $1 + md = 1 + k$, covolume $U$, and term count $1 + m(2^d - 1)$.
Setting $(m,d) = (1,k)$ recovers the pure hypercube; $(m,d) = (k,1)$ recovers
the pure star/chain.
\end{definition}

\begin{theorem}[Global optimality of the chain structure]
\label{thm:global}
For every $k \ge 1$ and every factorisation $k = md$ with $m, d \ge 1$,
\[
  g(m,d) \;:=\; \ln\bigl(1 + m(2^d - 1)\bigr) - \ln(1 + md) \;\ge\; 0,
\]
with equality if and only if $d = 1$.  Consequently,
$\alpha + \beta - 1 = g(m,d)/\ln n + o(1)$ is minimised uniquely at $d=1$ (the
pure star/chain) and maximised uniquely at $(m,d) = (1,k)$ (the pure hypercube).
\end{theorem}

\begin{proof}
It suffices to show $2^d - 1 \ge d$ for all positive integers $d$, with equality
only at $d = 1$.

\textit{Base:} $d=1$: $2^1 - 1 = 1 = d$. \checkmark

\textit{Inductive step:} Suppose $2^d - 1 \ge d$.  Then
$2^{d+1} - 1 = 2(2^d - 1) + 1 \ge 2d + 1 > d + 1$ for $d \ge 1$.

Hence $2^d - 1 \ge d$ with strict inequality for $d \ge 2$.  Since $\ln$ is
monotone, this gives $g(m,d) \ge 0$ (using $m(2^d-1) \ge md$) with equality
iff $d = 1$.

\textbf{Global maximum:}  For fixed $k = md$, $\ln(1+md) = \ln(1+k)$ is
constant, so maximising $g$ over factorisations of $k$ is equivalent to
maximising $f(d) := (2^d-1)/d$.  Since
\[
  \frac{f(d+1)}{f(d)}
  = \frac{d(2^{d+1}-1)}{(d+1)(2^d-1)} > 1
  \quad \text{for all } d \ge 1
\]
(as this simplifies to $2^d(d-1) > -1$, trivially true), $f$ is strictly
increasing, so $g$ is uniquely maximised at $d = k$ (forcing $m=1$), the pure
hypercube.  \qed
\end{proof}

\begin{corollary}
The pure hypercube construction is the unique \emph{worst} member of the
blocked hypercube--star family for every fixed total budget $k$, and every
interior point strictly interpolates between it and the chain's optimal value.
\end{corollary}

\begin{remark}[$U$-invariance and scope of Theorem~\ref{thm:global}]
The covolume $U$ is exactly the primorial of $n$ for \emph{every} configuration
$(m,d)$, every partition, and every seed.  This follows algebraically from the
lower-triangular basis structure and is not an assumption.  We verified this
computationally across 12 random partitions of varying $k$: $U$ equalled the
primorial in every case, even as the $W_i$ values differed by orders of
magnitude between seeds.

What Theorem~\ref{thm:global} does \emph{not} control is the \emph{achieved}
$\alpha$ for any specific partition, which depends on the lattice's first minimum
$\lambda_1$ relative to the Minkowski bound $C \cdot U^{1/(k+1)}$.  Different
partitions at fixed $(m,d)$ give different $W_i$ values and hence different
lattice geometries, with $\lambda_1$ varying accordingly --- this is the
Hermite-factor variance quantified in Section~\ref{sec:hermite}.  The theorem
compares proven \emph{upper bounds} on $\alpha$, not achieved values.
\end{remark}

\subsection{Numerical experiments and the Hermite factor}
\label{sec:hermite}

We implemented LLL lattice reduction using \texttt{mpmath} arbitrary-precision
arithmetic (with precision set to the lattice bit-length plus a safety margin of
$80$ bits), supplemented by a BKZ-style block enumeration pass over windows of
four consecutive basis vectors~\cite{schnorr1994}.  Plain \texttt{float64}
arithmetic was found to silently produce wrong results: at $576$-bit lattice
entries, all useful precision is lost under a $53$-bit mantissa, yielding
confident but meaningless output.

\paragraph{Results at $n = 400$.}  At $k = 19$ ($\beta \approx 0.500$, aiming
for the balanced point $\alpha = \beta = 1/2$), eight independent random
partition seeds gave implied exponent $1 + \max(\alpha,\beta)$ with mean $1.638$
and standard deviation $0.019$ (range $[1.602, 1.666]$).  At $k = 25$ (five
seeds), the mean was $1.612$ (sd $0.023$, range $[1.581, 1.665]$).

\paragraph{The Hermite-factor framing.}
The asymptotic convergence to exponent $3/2$ as $\beta \to 1/2$ is
\emph{proven unconditionally} by Theorem~\ref{thm:fullline} and does not depend
on these experiments.  The experiments measure our specific algorithm's
performance relative to the Minkowski-predicted magnitude $\log_2 U / (k+1)$.
At $k = 19$ this ratio was $2.305 \pm 0.255$; at $k = 25$ it was $2.573 \pm
0.347$.  The ratio did \emph{not} improve from $k = 19$ to $k = 25$: the
decrease in the raw exponent reflects the theorem's own shrinking target, not
the algorithm getting relatively closer.  Closing the Hermite-factor gap would
require a qualitatively stronger reduction algorithm (larger BKZ blocks, true
enumeration or sieving), not simply larger $k$.

%% ═══════════════════════════════════════════════════════════════════════════
\section{Seven Converging Obstructions}

Table~\ref{tab:summary} summarises the principal approaches.  Seven
structurally independent mathematical traditions all converge at or above
$\alpha + \beta \approx 1$, each for a distinct and individually
well-understood reason.

\begin{table}[ht]
\centering
\small
\begin{tabular}{@{}lllc@{}}
\toprule
Approach & Mathematical tradition & Achieved $\alpha+\beta$ & Verdict \\
\midrule
Naive bucketing          & Elementary number theory      & $= 1$ exactly    & Baseline \\
Lemma~6.2 CRT/CF         & CRT + continued fractions     & $\approx 1$, small savings & Best practical \\
Discrete-log sieve       & Fermat / multiplicative orders & $>1$ (worse)    & Ruled out \\
First-moment prob.\ method & Erd\H{o}s probabilistic method & $\to 1$       & Confirms barrier \\
AP-structured randomness & Structured probability        & $> 1$            & Ruled out \\
Dispersion (2nd moment)  & Linnik's method               & Wrong sign       & Mechanistically inapplicable \\
Hypercube lattice        & Minkowski / geometry of numbers & $=1$ exactly   & \textbf{Proved} \\
Star/chain               & Addition-chain theory         & $=1+o(1)$ for all $\beta$ & \textbf{Proved} \\
\bottomrule
\end{tabular}
\caption{Summary of approaches and their obstruction results.}
\label{tab:summary}
\end{table}

The paper's own combinatorial impossibility result (Lemma~6.3 of~\cite{umans2025})
and our probabilistic floor (Section~5) independently confirm the same barrier.
A direct test of the $c = 1$ edge case explicitly flagged as unresolved
in~\cite{umans2025} found the achieved offset never within an order of magnitude
of the necessary theoretical floor across all tested configurations, providing
empirical (not formal) evidence against that sub-case.

%% ═══════════════════════════════════════════════════════════════════════════
\section{Conclusion}

This investigation did not resolve the Umans--Wang $(\alpha,\beta)$-Divisor
Conjecture.  It produced three new proved results: Theorems~\ref{thm:hypercube},
\ref{thm:fullline}, and~\ref{thm:global}, together giving an exact
characterisation of, and global optimality result for, the blocked
hypercube--star family of lattice constructions.  Most substantively, it
produced a multi-method triangulation of a single barrier $\alpha + \beta \approx
1$, approached independently through combinatorics, probability theory,
Diophantine approximation, discrete logarithms, and the geometry of numbers,
with each approach failing for a distinct and well-understood reason.

\subsection{A calibrated self-assessment}
\label{sec:selfassess}

Theorems~\ref{thm:hypercube} and~\ref{thm:global} rest on elementary calculus
and a one-line inductive inequality; both were independently re-derived and
numerically verified.  Theorem~\ref{thm:fullline} is correct as a
pointwise-in-$\beta$ asymptotic statement, with the slowing convergence rate
near $\beta \to 1$ identified explicitly.  The non-degeneracy lemma
(Lemma~\ref{lem:nondegen}) is now fully proved by an elementary pigeonhole
argument, verified computationally and on a worked adversarial example; it rests
on one cited classical fact (Lemma~\ref{lem:bounded-basis}) rather than a fully
self-contained derivation.

The numerical results of Section~\ref{sec:hermite} characterise our specific
reduction algorithm's Hermite-factor gap to a proven asymptotic target; the
achieved-to-predicted ratio did not improve with $k$ over the tested range, so
these experiments should not be read as independent evidence of convergence to
exponent $3/2$.  The literature search supporting novelty claims was targeted,
not exhaustive.  Most importantly, this document has not been reviewed by anyone
outside the conversation in which it was produced.  Its appropriate use is as a
careful starting point for expert review, not as an established result.

%% ── References ────────────────────────────────────────────────────────────
\bibliographystyle{plainnat}
\begin{thebibliography}{99}

\bibitem{umans2025}
C.~Umans and J.~Wang.
\newblock A number-theoretic conjecture implying faster algorithms for
  polynomial factorization and integer factorization.
\newblock \emph{arXiv:2511.10851}, 2025.

\bibitem{kedlaya2011}
K.~Kedlaya and C.~Umans.
\newblock Fast polynomial factorization and modular composition.
\newblock \emph{SIAM J.\ Comput.}, 40(6):1767--1802, 2011.

\bibitem{knuth1969}
D.~Knuth.
\newblock \emph{The Art of Computer Programming, Volume~2: Seminumerical
  Algorithms}.
\newblock Addison-Wesley, 1969.

\bibitem{downey1981}
P.~Downey, B.~Leong, and R.~Sethi.
\newblock Computing sequences with addition chains.
\newblock \emph{SIAM J.\ Comput.}, 10(3):638--646, 1981.

\bibitem{linnik1963}
Y.~V.~Linnik.
\newblock \emph{The Dispersion Method in Binary Additive Problems}.
\newblock AMS, 1963.
\newblock (English translation of the 1961 Russian original.)

\bibitem{khinchin1964}
A.~Ya.~Khinchin.
\newblock \emph{Continued Fractions}.
\newblock University of Chicago Press, 1964.

\bibitem{postnikov1966}
A.~G.~Postnikov.
\newblock Ergodic problems in the theory of congruences and of Diophantine
  approximations.
\newblock \emph{Proc.\ Steklov Inst.\ Math.}, 82, 1966.

\bibitem{minkowski1896}
H.~Minkowski.
\newblock \emph{Geometrie der Zahlen}.
\newblock Teubner, 1896.

\bibitem{cassels1959}
J.~W.~S.~Cassels.
\newblock \emph{An Introduction to the Geometry of Numbers}.
\newblock Springer, 1959.

\bibitem{kz1873}
A.~Korkine and G.~Zolotareff.
\newblock Sur les formes quadratiques.
\newblock \emph{Math.\ Ann.}, 6:366--389, 1873.

\bibitem{viazovska2017}
M.~Viazovska.
\newblock The sphere packing problem in dimension~8.
\newblock \emph{Ann.\ Math.}, 185(3):991--1015, 2017.

\bibitem{erdos1947}
P.~Erd\H{o}s.
\newblock Some remarks on the theory of graphs.
\newblock \emph{Bull.\ Amer.\ Math.\ Soc.}, 53:292--294, 1947.

\bibitem{pollard1974}
J.~M.~Pollard.
\newblock Theorems on factorization and primality testing.
\newblock \emph{Math.\ Proc.\ Cambridge Philos.\ Soc.}, 76(3):521--528, 1974.

\bibitem{lll1982}
A.~K.~Lenstra, H.~W.~Lenstra, and L.~Lov\'asz.
\newblock Factoring polynomials with rational coefficients.
\newblock \emph{Math.\ Ann.}, 261:515--534, 1982.

\bibitem{schnorr1994}
C.~P.~Schnorr and M.~Euchner.
\newblock Lattice basis reduction: Improved practical algorithms and solving
  subset sum problems.
\newblock \emph{Math.\ Programming}, 66:181--199, 1994.

\end{thebibliography}

\end{document}
